% revised: CC 8 October 2026
% What this holds: Subsection Race and Ethnicity.
% Read by arlington-bsap.tex. Prose lives here; formatting lives in the wrapper and style/.
\subsection{Race and Ethnicity}

% The County's reply on the two lists (member-demographics-lists)
% gates the Board half of this section. Errors are corrected and the prose tightened;
% no claim is added or changed except where a source contradicted the draft.

\begin{figure}[t]
  \centering
  \refstepcounter{figure}\label{fig:residents-race}
  \figuretitle{County Population by Race and Ethnicity}
  \medskip
  \includegraphics[width=\textwidth]{residents_by_race.pdf}
  \smallskip
  \figurenotes{
    Notes: Panel (a) gives counts and panel (b) shares. Hispanic origin
    was first asked of every household in 1980, where the dashed rule
    falls; before then Hispanic residents were counted as Black or White,
    mostly White. White counts above 50,000 run off the top of panel (a);
    they peak at 161,329 in 1970.}{Source: US decennial census, 1870--2020.}
\end{figure}

Arlington was a majority-Black county through 1880 and half Black in 1890.
Between 1870 and 1950 the Black share of its residents fell from 63 percent to
4.8 percent, though the Black population itself grew from 2,010 to 6,517:
everything else grew faster. The share then barely moved, to 5.8 percent in
1970, and since 1980 the county has diversified mostly through Hispanic and
Asian American and Pacific Islander growth.

The county figures end where neighborhoods begin. A county count that rises
can hide a neighborhood cleared while in-migration elsewhere keeps the total
climbing. Bestebreurtje's
dissertation on African American community development in Arlington tells that
neighborhood history.\autocite{bestebreurtje2017}

% waits on: member-demographics-lists
% Generated by code/analysis/members_roster.py from data/clean/members.csv.
% Never edit by hand: bash run.sh rewrites it.
\begin{table}[t]
\refstepcounter{table}\label{tab:subsets}
\tabletitle{Women and Members of Color on the Board}
\medskip
\begingroup\footnotesize\setlength{\tabcolsep}{8pt}
\begin{tabular}{Q}
\toprule
\textit{Name} & \textit{Served} & \textit{Seated By} & \textit{Born} & \textit{Gender} & \textit{Race and Ethnicity} \\
\midrule
\multicolumn{6}{@{}l}{\bfseries 1800s} \\*
William A. Rowe & 1871--1883 & election & 1834\textsuperscript{c} & man\textsuperscript{c} & Black\textsuperscript{c,\,p} \\
John B. Syphax & 1872 & election & 1839\textsuperscript{c} & man\textsuperscript{c} & Black\textsuperscript{c,\,p} \\
Travis B. Pinn & 1879--1881 & election & 1841\textsuperscript{c} & man\textsuperscript{c} & Black\textsuperscript{c,\,p} \\
John W. Pendleton & 1883--1884 & election & 1848\textsuperscript{c} & man\textsuperscript{c} & Black\textsuperscript{c,\,p} \\
Tibbett Allen & 1887--1888 & election & 1840\textsuperscript{c} & man\textsuperscript{c} & Black\textsuperscript{c,\,p} \\
Walter G. Willson & 1889--1892 & election & --- & man\textsuperscript{a} & Black\textsuperscript{p} \\
\addlinespace[1.4ex]
\multicolumn{6}{@{}l}{\bfseries 1900s} \\*
Elizabeth B. Magruder & 1932--1947 & election & 1878\textsuperscript{c} & woman\textsuperscript{c,\,p} & White\textsuperscript{c} \\
Florence Cannon & 1948--1951 & election & 1889\textsuperscript{c} & woman\textsuperscript{c,\,p} & White\textsuperscript{c} \\
Leone B. Buchholz & 1952--1958 & special election & 1899\textsuperscript{c} & woman\textsuperscript{c,\,p} & White\textsuperscript{c} \\
Ellen Bozman & 1974--1997 & election & 1925\textsuperscript{p} & woman\textsuperscript{p} & White\textsuperscript{a} \\
Dorothy Grotos & 1976--1983 & election & 1931\textsuperscript{p} & woman\textsuperscript{p} & White\textsuperscript{a} \\
Mary Margaret Whipple & 1983--1995 & election & 1941\textsuperscript{p} & woman\textsuperscript{p} & White\textsuperscript{a} \\
William T. Newman, Jr. & 1988--1993 & election & 1950\textsuperscript{p} & man\textsuperscript{p} & Black\textsuperscript{p} \\
Barbara A. Favola & 1997--2011 & special election & 1955\textsuperscript{p} & woman\textsuperscript{p} & White\textsuperscript{a} \\
\addlinespace[1.4ex]
\multicolumn{6}{@{}l}{\bfseries 2000s} \\*
Charles P. Monroe & 2000--2003 & election & 1957\textsuperscript{p} & man\textsuperscript{p} & Black\textsuperscript{p} \\
J. Walter Tejada & 2003--2015 & special election & 1958\textsuperscript{p} & man\textsuperscript{p} & Hispanic\textsuperscript{p} \\
Mary H. Hynes & 2008--2015 & election & 1956\textsuperscript{p} & woman\textsuperscript{p} & White\textsuperscript{a} \\
Libby T. Garvey & 2012--2024 & special election & 1951\textsuperscript{p} & woman\textsuperscript{p} & White\textsuperscript{a} \\
Christian E. Dorsey & 2016--2023 & election & 1972\textsuperscript{p} & man\textsuperscript{p} & Black\textsuperscript{p} \\
Kate A. "Katie" Cristol & 2016--2023 & election & 1985\textsuperscript{p} & woman\textsuperscript{p} & White\textsuperscript{a} \\
Tannia Talento & 2023 & appointment & --- & woman\textsuperscript{p} & Hispanic\textsuperscript{p} \\
Maureen E. Coffey & 2024-- & election & 1995\textsuperscript{p} & woman\textsuperscript{p} & White\textsuperscript{a} \\
Susan R. Cunningham & 2024-- & election & --- & woman\textsuperscript{p} & White\textsuperscript{a} \\
Julius D. "JD" Spain, Sr. & 2025-- & election & 1972\textsuperscript{p} & man\textsuperscript{p} & Black\textsuperscript{p} \\
\bottomrule
\end{tabular}\endgroup
\par\smallskip
\figurenotes{Notes: Superscripts indicate the source for each record: the US census (c), press or published profile (p), or assumption (a). Two superscripts mean two kinds of source agree. A dash means no source gives the value. This table is a subset of the full roster, Table~\ref{tab:roster} in the Data Appendix.}{}
\end{table}


For the Board's first 14 years Black members held one or two of its three seats,
though never in proportion: from 1871 to 1892 they held, year by year, about
30 percent of the Board's seats in a county whose Black share averaged about 55
percent. Walter G.
Willson, who won the Washington district seat in 1889 and 1891 and died in
office in June 1892, was the last. No source names another member of color until
William T. Newman, Jr. in 1988, 96 years later. Since then no more than one has
sat at a time except in the second half of 2023, when Dorsey and Talento
overlapped, and none sat in 2024. Figure~\ref{fig:board-race-gender} counts every
member not identified as a member of color as White, an assumption the Data
Appendix explains.

The break came in two steps, and the sources explain neither. In 1888 the county
court issued a rule against Tibbett Allen, the Board's only Black member, over
his residence; he resigned, and the judge appointed Frank Hume to Jefferson
district.\autocite{alexandriagazette1888rule} Hume then won the district in 1889, 1891 and 1893,
with no opponent recorded in 1889 or 1893.\autocite{alexandriagazette18890525p3,%
alexandriagazette18930526p3} Jefferson, 65 percent Black in 1880, had elected a
Black member at 10 of its 11 elections through 1887, and Allen's removal ended
that. Willson held a seat in Washington district until he died, and nothing
consulted says why no Black candidate won a seat after him.

Four forces may have kept the Board white once Willson was gone. (1) In 1932
the county moved to at-large election, which removed the district through which
a concentrated minority could elect a member. (2) The county's Black population
fell, from a majority in 1870 to one resident in eight by 1930. (3) The
electorate shrank after the Walton Act of 1894 and the constitution of 1902,
taking an unknown number of Black voters with it. (4) The three magisterial
districts never guaranteed a Black seat, even where Black residents were the
majority.

At-large election in 1932 cannot be tested as a cause of exclusion, because
exclusion was complete before it came. By 1931 no Black member had sat for 39
years, and no source names a Black candidate after 1903. The three who ran in
November 1931 (Mary B. Harris, E. T. Morton and C. H. Moseley) each finished
below sixth place among 51 candidates for five seats, in a county 12.5 percent
Black.\autocite[67]{anderson1958} What 1932 changed was the route: at-large
election removed the district through which a concentrated minority could elect
a member. Staggering, which county voters approved in 1938 by 1,539 to 1,487 in
the county's tally, then raised the threshold, because most elections filled a
single seat and each needed a countywide
majority.\autocite[11]{arlingtonelections2021} No Black candidate ran for the
Board from 1932 until 1987, so nothing shows what either change would have done
to one.\autocite[22--23]{pratt1995}\fnsep\autocite[217--218]{bestebreurtje2017}
George Vollin, testifying in 1974 in the vote-dilution suit he brought,
recalled that after the defeats of 1931 Black Arlingtonians considered a
countywide race an exercise in futility, since Whites greatly outnumbered
them.\autocite[22--23]{pratt1995}

Population change was the background, not the break. The county's Black share
fell from 63 percent in 1870 to 12.5 percent in 1930, but \blackShareMenJeffersonNineteenHundred{}
percent of the men of voting age in Jefferson district were Black in 1900
(Figure~\ref{fig:residents-district-race}), and in the eight Jefferson
contests from 1889 through 1903 no Black candidate won. A district that was
still majority Black sent no Black member to the Board.

\begin{figure}
  \centering
  \refstepcounter{figure}\label{fig:residents-district-race}
  \figuretitle{Black Share of Men of Voting Age by Magisterial District}
  \medskip
  \includegraphics[width=\textwidth]{residents_by_district_race_adults.pdf}
  \smallskip
  \figurenotes{
    Notes: Each line divides a district's Black men aged 21 and over by all of
    its men aged 21 and over. A dotted segment crosses a census that gives no
    district count: the 1890
    schedules burned, and Arlington district's 1900 schedules are short by about
    one resident in six. The open rings at 1870 are the Black share of all
    residents of each township, of every age and sex, not of men aged 21 and
    over, so each line's first segment joins two different counts: the 1870 volume prints race by township for residents only, and the
    census file places no one in a township. The county line is the three districts added together,
    and Alexandria's annexations of 1915 and 1930 mean no district covers the
    same land throughout. Districts ended in 1932, and the 1930 census has no
    district split by race.}{Source: US decennial census, 1870--1920.}
\end{figure}

The electorate collapsed, and the laws that shrank it fit the years the absence
lasted. In 1880 the county cast votes for President equal to
\turnoutEighteenEighty{} percent of its men of voting age
(Figure~\ref{fig:turnout-president}); by 1892 the figure was
\turnoutEighteenNinetyTwo{} percent, in 1900 \turnoutNineteenHundred{} percent,
and in 1904 \turnoutNineteenFour{} percent, where it stayed, at
\turnoutNineteenEight{} and \turnoutNineteenTwelve{} percent, through 1912.
The General Assembly's Walton Act of 1894 made the
official ballot the only legal one. It required a voter to strike out the name
of every candidate he did not want with a line three-fourths of the name's
length, voided any office with more than one name left standing, and let only a
constable the electoral board appointed help a voter who could not read the
ballot.\autocite[secs.~3, 11, 15]{waltonact1894} The constitution of 1902
required registration and a state poll tax paid six months before the election,
and, from 1904, an application written in the applicant's own hand ``without
aid, suggestion, or memorandum.''\autocite[art.~II, secs.~18--21]{vaconstitution1902}
Neither law can explain 1888, which came six and 14 years before them. No source
consulted counts how many of the men who stopped voting were Black, but Jefferson
district, the county's most heavily Black, shows the fall most steeply. Its
1901 Board contest drew votes equal to \boardTurnoutJeffersonNineteenOne{}
percent of the district's men of voting age; its 1907 contest drew
\boardTurnoutJeffersonNineteenSeven{} percent, and the seat turned on 2 votes
of 161. In a district that had sent
Black members to the Board at nearly every election a generation earlier, a
handful of men voting together could now decide it.

The figure opens high because the country's electorate was. The Hayes-Tilden
election of 1876 drew about 83 percent of eligible voters nationally, the
highest turnout in American history, with Black men voting under federal
protection and both parties organized to the last voter.\autocite{burnham2016turnout}
The county's \turnoutEighteenSeventySix{} votes per 100 men that year
is that electorate, before the laws of 1894 and 1902 took it apart.

\begin{figure}
  \centering
  \refstepcounter{figure}\label{fig:turnout-president}
  \figuretitle{Arlington Voter Turnout in Presidential Elections, 1872--2024}
  \medskip
  \includegraphics[width=\textwidth]{elections_turnout_president.pdf}
  \smallskip
  \figurenotes{Notes: This graph plots votes cast for President per 100 residents of
  voting age in Arlington County. Residents of voting age are drawn from US census data, with a straight
  line between censuses. They include men aged 21 and older through 1916; all
  residents 21 and older from 1920, following the 19th Amendment; and all
  residents 18 and older from 1971, following the 26th Amendment. They include
  residents who were ineligible to vote because of citizenship status, criminal
  history and other legal restrictions; the census does not record those.
  Election data are detailed in the appendix. Dashed rules mark the Walton Act of 1894, the constitution
  of 1902 and the first election with women voting. On the same denominator
  1876 reads 76 for Virginia, 77 to 83 for Fairfax,
  Loudoun and Prince William, and 91 for Alexandria city.}{Source:
  Warrock-Richardson Almanack; Secretary of the Commonwealth; Virginia
  Department of Elections; \textit{Alexandria Gazette}; O'Leary (2010);
  Arlington County, \textit{Arlington Elections}; US decennial census,
  1870--2020.}
\end{figure}

The districts never guaranteed a seat either. Jefferson elected a Black member at 10
of its 11 elections through 1887, but Arlington district, 62 and 56 percent
Black in 1870 and 1880, did so at 3 of 12, and Washington, the least Black of
the three, at none until Willson. In 1885 William A. Rowe lost Arlington
district, as the regular Republican nominee, to a ``citizens' candidate.''%
\autocite{alexandriagazette18850529p3} A district gave a concentrated minority a
route to a seat, and even a district with a Black majority could go the other
way.

What moved county voters to choose at-large election in 1930 is known only from
testimony. Pratt writes that countywide election substantially reduced Black
voting strength compared with a district of segregated neighborhoods, but his
account of the motive is the 1974 testimony of George Vollin and Harrison
Douglas, recalled more than 40 years after the vote, and he adds that whether
the change caused the defeat of 1931 ``may never be conclusively
known.''\autocite[22]{pratt1995} The court ruling Cornelia Rose reads into the
choice is in Section~\ref{sec:method}.

% waits on: school-board-elections-1947-1956

The two Vollin actions are separate cases. The 1974 federal suit brought the
vote-dilution claim, and it failed. The 1976 state case asked whether voters
could use the referendum statute a second time to put district election to a
vote, and the courts said no.

The sources settle three things. The absence began in 1888, before either
disfranchising law and in a district that was still a Black majority; neither
the fall in the county's Black population nor the district system, taken alone,
explains it; and the collapse of the electorate after 1902 is the one force
that fits how long it lasted. That is a fit and not a count, because
registration and poll-tax records by race, which would show it, are in no source
consulted. Nor does any source show whether Black candidates stopped running
between 1892 and 1930 or ran and lost, or why Allen was removed.

The Historical Society dates Newman's election 1987, Monroe's 1999, Dorsey's
2015 and Spain's 2024, and this study lists each as taking the seat a year later. Both
are correct: each won in November and took office the following January.

